\documentclass[11pt,a4paper]{article}

\usepackage[margin=2.4cm]{geometry}
\usepackage{amsmath,amssymb}
\usepackage{booktabs}
\usepackage{longtable}
\usepackage{array}
\usepackage{xcolor}
\usepackage{tcolorbox}
\usepackage[colorlinks=true,linkcolor=black,citecolor=black,urlcolor=blue]{hyperref}
\usepackage{fancyhdr}
\usepackage{caption}
\usepackage{enumitem}

\definecolor{paafblue}{HTML}{2563EB}
\definecolor{paafslate}{HTML}{0F1B2D}
\definecolor{paafgrey}{HTML}{64748B}
\definecolor{warnbg}{HTML}{FEF3C7}
\definecolor{warnbd}{HTML}{D97706}

\newtcolorbox{caveat}[1][]{colback=warnbg!40,colframe=warnbd,
  boxrule=0.6pt,arc=2pt,left=6pt,right=6pt,top=5pt,bottom=5pt,#1}

\newtcolorbox{algobox}[1][]{colback=white,colframe=paafslate,
  boxrule=0.6pt,arc=2pt,left=8pt,right=8pt,top=6pt,bottom=6pt,#1}

% Pseudocode without algorithm.sty: a stepped list with its own counter.
\newcounter{algoline}
\newcommand{\algoreset}{\setcounter{algoline}{0}}
\newcommand{\alg}[2][0]{%
  \refstepcounter{algoline}%
  \noindent\makebox[2.2em][r]{\scriptsize\color{paafgrey}\thealgoline:}%
  \hspace{4pt}\hspace{#1\parindent}\parbox[t]{\dimexpr\linewidth-2.2em-4pt-#1\parindent}{\raggedright #2}\par\vspace{1.2pt}}
\newcommand{\kw}[1]{\textbf{\color{paafblue}#1}}

\pagestyle{fancy}
\fancyhf{}
\lhead{\small\color{paafgrey}PAAF --- Amorphous Cell Construction}
\rhead{\small\color{paafgrey}\thepage}
\renewcommand{\headrulewidth}{0.3pt}

\captionsetup{font=small,labelfont=bf}
\setlist{itemsep=2pt,parsep=0pt,topsep=3pt}

\title{\vspace{-1.2cm}\textbf{Amorphous Cell Construction in PAAF}\\[4pt]
\large Theory, Algorithm, Implementation and Validation}
\author{Polymer Auto-Assembly Framework (PAAF)}
\date{\today}

\begin{document}
\maketitle
\thispagestyle{fancy}

\begin{abstract}
\noindent
PAAF builds periodic amorphous polymer cells by \emph{chain growth}: each
chain is constructed one skeletal bond at a time, with the torsion about that
bond drawn from a rotational isomeric state (RIS) distribution reweighted by
the non-bonded environment already present in the box. This is the method of
Theodorou and Suter~\cite{ts1985}, optionally combined with the scanning
method of Meirovitch~\cite{meirovitch1983}, on the statistical foundation laid
by Flory~\cite{flory1969}. It is the same published pair of algorithms that
underlies the BIOVIA Materials Studio \emph{Amorphous Cell} module, which
BIOVIA documents as ``the combined use of the arc algorithm developed by
Theodorou and Suter and the scanning method of Meirovitch''~\cite{biovia}.

This report states the theory with its equations, specifies exactly what PAAF
implements, reports the measured validation of every physical claim, and sets
out --- without softening --- what the implementation does \emph{not} do. The
characteristic ratio $C_\infty$ is reproduced within $12\%$ of tabulated
values for all four parameterised polymers; bond lengths and angles are exact
to floating-point precision; melt densities are reached, unbiased at
$0.85$ and $1.04$~g/cm$^3$. The principal limitations are that PAAF performs \emph{construction
only} with no subsequent energy minimisation, that growth is at
united-skeletal-atom resolution followed by back-mapping rather than direct
full-atomistic construction, and that polymer-specific RIS parameters exist
for four polymers only. Section~\ref{sec:novelty} assesses novelty honestly:
the method is a faithful reimplementation, not a new algorithm.
\end{abstract}

\tableofcontents
\vspace{6pt}
\hrule

%======================================================================
\section{Why growth rather than insertion}
\label{sec:why}

The natural first approach to filling a periodic box with polymer is to build
whole chains and insert them at random positions and orientations, rejecting
overlaps. This works for solvents and short oligomers and fails for polymers,
for a reason that is geometric rather than a matter of tuning.

A random coil pervades a volume set by its radius of gyration,
$R_g \propto N^{1/2}$, so the pervaded volume scales as $N^{3/2}$ while the
mass inside it scales only as $N$. At melt density that pervaded volume is
mostly \emph{other chains}: a polymer coil is roughly $99\%$ not-itself.
Inserting a chain whole therefore requires locating a cavity shaped like an
entire coil, and the probability of that decreases roughly exponentially in
$N$. Acceptance collapses somewhere between $20$ and $50$ repeat units.

Growing the chain bond by bond replaces one impossible question --- ``is there
a coil-shaped hole?'' --- with many easy ones --- ``is there room for one more
skeletal atom?''. Each step requires only local free volume, and the chain
finds its own way through the material already placed. This is the central
insight of Theodorou and Suter~\cite{ts1985}, and it is why every serious
amorphous cell builder grows rather than inserts.

%======================================================================
\section{Rotational isomeric state theory}
\label{sec:ris}

\subsection{The three-state model}

A backbone torsion is not free. Rotation about a C--C bond has three minima
--- \emph{trans} ($\phi \approx 180^\circ$) and \emph{gauche}$^{\pm}$
($\phi \approx \pm 60^\circ$) --- separated by barriers of a few kcal/mol. The
RIS approximation of Flory and Volkenstein replaces the continuous torsional
potential by these discrete states~\cite{flory1969,mattice1994}. Two facts
follow, and both control chain dimensions.

First, the minima are not equally weighted. A gauche state costs an extra
energy $E_\sigma$ relative to trans, giving the first-order statistical weight
\begin{equation}
  \sigma = \exp\!\left(-\frac{E_\sigma}{RT}\right) < 1 .
  \label{eq:sigma}
\end{equation}

Second, neighbouring torsions are \emph{correlated}. A $g^{+}g^{-}$ pair
drives the two chain segments on either side into one another --- the
\emph{pentane effect} --- and carries a large additional penalty $E_\omega$,
\begin{equation}
  \omega = \exp\!\left(-\frac{E_\omega}{RT}\right) \ll 1 .
  \label{eq:omega}
\end{equation}

Both are encoded in the $3\times3$ statistical weight matrix, indexed
$U[\,\text{previous state},\ \text{current state}\,]$ in the order
$(t,\, g^{+},\, g^{-})$:
\begin{equation}
  \mathbf{U}(T) =
  \begin{pmatrix}
    1 & \sigma & \sigma \\
    1 & \sigma & \sigma\omega \\
    1 & \sigma\omega & \sigma
  \end{pmatrix}.
  \label{eq:umatrix}
\end{equation}
Every row begins with $1$ because trans is the reference state. The factor
$\omega$ appears only on the $g^{+}g^{-}$ and $g^{-}g^{+}$ entries --- exactly
the pairs whose substituents collide. In PAAF at $T = 413$~K for
polyethylene this gives $U[g^{+},g^{+}] = 0.5438$ against
$U[g^{+},g^{-}] = 0.0421$, a suppression of $12.9\times$, which is the pentane
effect made numerical.

\subsection{Exact finite-chain sampling}

A naive sampler would draw the next state from $P(k) \propto U[j,k]$. That is
wrong for a chain of finite length, because it ignores how many conformations
each choice leaves available downstream, and it biases the chain ends.

PAAF instead uses Flory's matrix method as a \emph{sampler}. The partition
function is accumulated backwards from the far end of the chain,
\begin{equation}
  \mathbf{Z}^{(i)} = \mathbf{U}\,\mathbf{Z}^{(i+1)},
  \qquad \mathbf{Z}^{(n)} = \mathbf{1},
  \label{eq:backward}
\end{equation}
after which the forward conditional probability is exact:
\begin{equation}
  P(k \mid j) =
  \frac{U[j,k]\, Z^{(i+1)}_{k}}
       {\sum_{k'} U[j,k']\, Z^{(i+1)}_{k'}} .
  \label{eq:conditional}
\end{equation}
The recursion in Eq.~\eqref{eq:backward} is renormalised at each step, since
the product underflows for long chains and only ratios are ever used.

Table~\ref{tab:pops} shows that the realised state populations from
Eq.~\eqref{eq:conditional} reproduce the transfer-matrix prediction --- the
exact long-chain limit obtained from the principal eigenvector of
$\mathbf{U}$ --- to better than $0.006$ in every state.

\begin{table}[t]
\centering
\small
\begin{tabular}{llrrr}
\toprule
Polymer & State & Realised & Predicted & $\Delta$ \\
\midrule
PE & $t$      & 0.6011 & 0.5974 & $+0.0037$ \\
   & $g^{+}$  & 0.2036 & 0.2013 & $+0.0023$ \\
   & $g^{-}$  & 0.1953 & 0.2013 & $-0.0060$ \\
\midrule
PS & $t$      & 0.6786 & 0.6741 & $+0.0045$ \\
   & $g^{+}$  & 0.1629 & 0.1630 & $-0.0001$ \\
   & $g^{-}$  & 0.1585 & 0.1630 & $-0.0044$ \\
\bottomrule
\end{tabular}
\caption{Sampled torsion-state populations against the transfer-matrix
prediction, $T = 413$~K, 80 chains of 500 torsions. Measured by
\texttt{tests/test\_ris.py}.}
\label{tab:pops}
\end{table}

\subsection{Chain geometry}

Bond lengths $l$ and bond angles $\theta$ are held \emph{fixed}; only the
torsion varies. Each successive skeletal atom is placed by exact internal
coordinates using the Natural Extension Reference Frame construction: given
three consecutive placed atoms $\mathbf{a}, \mathbf{b}, \mathbf{c}$, the next
atom is
\begin{equation}
  \mathbf{d} = \mathbf{c} + l\,\mathbf{R}(\mathbf{a},\mathbf{b},\mathbf{c})
  \begin{pmatrix}
    -\cos\theta \\ \sin\theta\cos\phi \\ \sin\theta\sin\phi
  \end{pmatrix},
  \label{eq:nerf}
\end{equation}
where $\mathbf{R}$ is the orthonormal frame built from the preceding three
atoms. Because Eq.~\eqref{eq:nerf} is an exact construction rather than a
minimisation, the geometry cannot drift: PAAF measures bond lengths of
$1.530000 \pm 1.5\times10^{-15}$~\AA{} and bond angles of
$112.0000 \pm 5.7\times10^{-14}\,^\circ$ over chains of 300 torsions, which
is floating-point noise.

\subsection{The characteristic ratio, and why it is the test}

The observable that reveals whether an RIS implementation is real is the
characteristic ratio
\begin{equation}
  C_n = \frac{\langle R^2 \rangle}{n\,l^2},
  \label{eq:cn}
\end{equation}
the mean-square end-to-end distance relative to a freely jointed walk of the
same bond count. Its limiting value $C_\infty$ is tabulated experimentally for
real polymers. Three regimes must be distinguishable:

\begin{center}
\small
\begin{tabular}{lcl}
\toprule
Model & $C$ for PE & What it captures \\
\midrule
Freely jointed & $1.0$ & nothing \\
Freely rotating ($\theta$ fixed) & $\approx 2.2$ & bond angle only \\
RIS ($\theta$ fixed, $\phi$ weighted) & $\approx 6.7$ & + torsional statistics \\
\bottomrule
\end{tabular}
\end{center}

The freely rotating limit has a closed form that serves as an independent
check on the geometry code alone:
\begin{equation}
  C_\infty^{\text{fr}} = \frac{1 + \cos\alpha}{1 - \cos\alpha},
  \qquad \alpha = 180^\circ - \theta .
  \label{eq:freelyrotating}
\end{equation}
For $\theta = 112^\circ$ this gives $2.198$; PAAF measures $2.445 \pm 0.098$
with torsions sampled uniformly, confirming that the chain builder is correct
before any statistics are applied. With RIS weights switched on, the same code
gives $6.82$ against $2.22$ for random torsions --- a factor of $3.07$. An
implementation in which RIS were decorative would sit at the freely rotating
baseline.

%======================================================================
\section{Growth into the periodic cell}
\label{sec:growth}

\subsection{The growth probability}

The Theodorou--Suter construction adds one skeletal bond at a time. For each
candidate rotational isomeric state $s$ the trial atom position is computed
from Eq.~\eqref{eq:nerf}, and the state is chosen with probability
\begin{equation}
  P(\phi_s) \;\propto\; U[\,j, s\,]\;
  \exp\!\left(-\frac{E^{\text{nb}}_{s}}{k_B T}\right),
  \label{eq:growth}
\end{equation}
where $j$ is the state of the previous torsion and $E^{\text{nb}}_{s}$ is the
non-bonded energy of the trial atom against the atoms of \emph{previously
completed chains}, evaluated under the minimum-image convention within the
neighbour-list cutoff. Equation~\eqref{eq:growth} is the heart of the method:
the first factor makes the coil the right \emph{size}, the second makes it
avoid \emph{other chains}. Note that the growing chain's own beads do not
enter $E^{\text{nb}}$ --- a chain is inserted into the neighbour grid only
once complete --- so intra-chain crowding is handled entirely by the hard
self-avoidance test described below.

\subsection{The non-bonded term}

A full Lennard-Jones sum would be the faithful choice, but it requires
per-atom LJ parameters that are not available until the force field is
assigned --- which in PAAF happens \emph{after} the cell is built. PAAF
therefore uses a purely repulsive soft-core term keyed on van der Waals
contact distance,
\begin{equation}
  E^{\text{nb}} = \epsilon \sum_{j}
  \left(\frac{\sigma_{ij}}{r_{ij}}\right)^{12},
  \qquad r_{ij} < \sigma_{ij},
  \label{eq:softcore}
\end{equation}
with $\epsilon = 0.2$~kcal/mol.

\begin{caveat}
\textbf{Deviation from the paper, stated plainly.}
Equation~\eqref{eq:softcore} reproduces the qualitative effect that matters
for construction --- candidates that bury the new atom in existing material
are strongly disfavoured --- without claiming a quantitative energy. Absolute
values of $E^{\text{nb}}$ are not physical energies and are not reported as
such anywhere in PAAF.
\end{caveat}

\subsection{Hard rejection: overlap and spearing}

Two hard tests are applied before Eq.~\eqref{eq:growth} is evaluated, because
no statistical weight however favourable should be able to produce them.

\paragraph{Overlap.} A candidate is rejected if it comes closer to any
non-bonded atom than
\begin{equation}
  r_{\min} = \max\!\left(\,d_{\text{tol}},\;
  \tfrac{1}{2}\,(R_i + R_j)\,s\,\right),
  \label{eq:overlap}
\end{equation}
with interpenetration factor $s = 0.55$ and hard tolerance
$d_{\text{tol}} = 1.7$~\AA{} by default. Neighbour searching against other
chains uses a cell list, so that part of the cost is $O(N)$ rather than
$O(N^2)$; the self-avoidance scan along the growing chain is not, and is
$O(n)$ per candidate. The 1--2, 1--3 \emph{and} 1--4 partners along the chain
are excluded --- the first two are at bonded distance by construction, and the
1--4 separation is set by the very torsion being chosen, so a gauche state
legitimately brings it inside the tolerance. Against its own chain the test
uses $d_{\text{tol}}$ alone; the radius term in Eq.~\eqref{eq:overlap} applies
to inter-chain contacts, and only exceeds $1.7$~\AA{} above roughly
$226$~amu per bead, which none of the polymers in Table~\ref{tab:gran} reach.

\paragraph{Ring spearing.} A threaded ring is a topological defect: the bond
cannot be pulled back out, so energy minimisation relaxes the strain and
leaves the two permanently interlocked. It is a classic silent failure of
constructed cells. Each ring is approximated by the disc spanned by its atoms
and the segment--disc intersection solved directly,
\begin{equation}
  t = \frac{\mathbf{n}\cdot(\mathbf{c}_0 - \mathbf{p}_0)}
           {\mathbf{n}\cdot(\mathbf{p}_1 - \mathbf{p}_0)},
  \qquad
  \text{speared} \iff
  0 \le t \le 1 \ \wedge\
  \lVert \mathbf{h} - \mathbf{c}_0 \rVert \le R_{\text{ring}},
  \label{eq:spear}
\end{equation}
where $\mathbf{h}$ is the crossing point.

\begin{caveat}
\textbf{Where this test can and cannot run.} At skeletal-bead resolution a
pendant ring is not represented at all --- polystyrene's phenyl group does not
exist until back-mapping places it --- so during growth there is literally
nothing to spear and the predicate, though implemented and unit-tested against
four geometric cases, never fires. PAAF therefore performs the check on the
finished \emph{all-atom} cell instead (Section~\ref{sec:backmap}), where it is
a diagnostic rather than a rejection criterion. This is worth doing: applying
it to a polystyrene cell at $1.04$~g/cm$^3$ immediately found \textbf{3
speared bonds} among 20 phenyl rings --- a defect the construction was
previously blind to. A full-atomistic grower could reject these during growth;
PAAF can only report them.
\end{caveat}

\begin{caveat}
\textbf{On the overlap tolerance.} $d_{\text{tol}}$ is \emph{not} the physical
contact distance and must not be set to one. Equation~\eqref{eq:softcore}
already discourages close approach; $d_{\text{tol}}$ is only a guard against
an overlap so severe that no subsequent minimisation could repair it. Setting
it to the true united-atom contact distance double-counts the repulsion and
prevents the cell being built at all. The default was measured, not chosen:
see Table~\ref{tab:tol}.
\end{caveat}

\begin{table}[t]
\centering
\small
\begin{tabular}{ccc}
\toprule
$d_{\text{tol}}$ (\AA) & Cells built & Mean $C_n$ \\
\midrule
2.0 & 1 / 3 & 5.68 \\
1.9 & 2 / 3 & 6.69 \\
1.8 & 2 / 3 & 7.60 \\
\textbf{1.7} & \textbf{3 / 3} & \textbf{5.74} \\
1.5 & 3 / 3 & 5.27 \\
\bottomrule
\end{tabular}
\caption{Effect of the hard overlap tolerance. Eight polyethylene chains of
$\mathrm{DP}=60$ grown into a $29.7$~\AA{} box at $0.85$~g/cm$^3$ (the melt
density at 413~K), unbiased sampling, three random seeds each. Increasing
\texttt{max\_restarts} instead is a worse remedy: at $2.0$~\AA{}, raising
restarts from 30 to 150 still reached only 2/3 and pulled $C_n$ down to
$4.97$, because surviving a restart selects for compact chains.}
\label{tab:tol}
\end{table}

\subsection{Meirovitch scanning and the Rosenbluth bias}

Growth can walk into a dead end from which no continuation is acceptable, and
the whole chain is then lost --- \emph{attrition}. The scanning method of
Meirovitch~\cite{meirovitch1983} looks $k$ steps ahead and weights each
candidate by the total surviving weight of its continuations,
\begin{equation}
  W_s^{(k)} = \sum_{\text{continuations } c}
  \prod_{m=1}^{k} U[\,s_{m-1}, s_m\,]\,
  \exp\!\left(-\frac{E^{\text{nb}}_{m}}{k_B T}\right),
  \label{eq:scan}
\end{equation}
so that Eq.~\eqref{eq:growth} is replaced by
$P(\phi_s) \propto U[j,s]\,\exp(-E^{\text{nb}}_s/k_BT)\,W_s^{(k)}$. The cost
is $3^{k}$ per step, so $k \in \{1,2,3\}$ is the practical range.

Scanning demonstrably works against attrition. Four polyethylene chains of
$\mathrm{DP}=40$ at $1.10$~g/cm$^3$ over three seeds: $k=0$ built 2 of 3 cells
with 23 restarts; $k=2$ built 3 of 3 with 3 restarts.

\begin{caveat}
\textbf{Scanning is off by default, and the reason is measured.}
Selecting a state in proportion to its look-ahead weight is a \emph{biased}
estimator --- the Rosenbluth bias --- because states that leave more room are
over-chosen, and those are the extended ones. Measured on dilute polyethylene
(2 chains of $\mathrm{DP}=40$ at $0.10$~g/cm$^3$, 12 seeds), against
$C_n = 6.78 \pm 0.24$ for the same chain grown unperturbed:
$5.66 \pm 0.74$ at $k=0$, $7.40 \pm 1.07$ at $k=1$, $7.52 \pm 1.35$ at $k=2$.
Scanning overshoots the unperturbed dimensions; $k=0$ undershoots them.
Removing the bias requires carrying the Rosenbluth weight
\begin{equation*}
  W = \prod_{i} \sum_{s} U[j_i,s]\,\exp(-E^{\text{nb}}_{i,s}/k_BT)
\end{equation*}
through the construction and using it in every subsequent average, which is
\emph{not} implemented in PAAF. Therefore: use $k=0$ when chain dimensions
matter; use $k>0$ only when attrition would otherwise prevent the cell being
built at all, and do not quote $C_n$ from such a cell. PAAF attaches this
warning to any result produced with $k>0$.
\end{caveat}

\subsection{Growth granularity: one bead per skeletal atom}

The RIS model parameterises a \emph{skeletal bond}. A growth step must
therefore be a skeletal bond, not a repeat unit. PAAF determines the number of
skeletal atoms a repeat unit contributes by finding the shortest path between
the two polymerisation attachment points in the SMILES graph:
\begin{equation}
  n_{\text{skel}} = \bigl|\,\mathrm{path}(a_1 \to a_2)\,\bigr| ,
  \qquad a_k = \text{heavy atom bonded to } \ast_k ,
  \qquad
  m_{\text{bead}} = \frac{M_{\text{unit}}}{n_{\text{skel}}} .
  \label{eq:granularity}
\end{equation}
Side groups are branches off that path and do not count.
Table~\ref{tab:gran} gives the resulting values.

\begin{caveat}
\textbf{This was a real defect, and it is worth recording why.} An earlier
version of PAAF placed one bead per \emph{repeat unit} at the RIS bond length
$1.53$~\AA. Density came out correct, because density knows only about mass
and volume. But the contour length was too short by a factor of
$n_{\text{skel}}$, so every dimension derived from the chain path --- $R_g$,
end-to-end distance, $C_n$ --- was too small by $\sqrt{n_{\text{skel}}}$:
$\sqrt{2}$ for the vinyl polymers, $\sqrt{10}$ for poly(butylene succinate).
A correct-looking density concealed it. The lesson generalises: check an
observable that depends on the chain \emph{path}, not only one that depends on
its mass.
\end{caveat}

\begin{table}[t]
\centering
\small
\begin{tabular}{llrrr}
\toprule
Polymer & Repeat unit SMILES & $n_{\text{skel}}$ & $M_{\text{unit}}$ & $m_{\text{bead}}$ \\
\midrule
Polyethylene            & \texttt{[*]CC[*]}                     & 2  & 28.05  & 14.03 \\
Polypropylene           & \texttt{[*]CC([*])C}                  & 2  & 42.08  & 21.04 \\
Polystyrene             & \texttt{[*]CC([*])c1ccccc1}           & 2  & 104.15 & 52.08 \\
Poly(vinyl chloride)    & \texttt{[*]CC([*])Cl}                 & 2  & 62.50  & 31.25 \\
Poly(ethylene oxide)    & \texttt{[*]OCC[*]}                    & 3  & 44.05  & 14.68 \\
Poly(butylene succinate)& \texttt{[*]OCCCCOC(=O)CCC(=O)[*]}     & 10 & 172.18 & 17.22 \\
\bottomrule
\end{tabular}
\caption{Skeletal granularity and masses derived from SMILES, in g/mol. Note
that the phenyl ring of polystyrene is a branch and correctly contributes
$0$ to $n_{\text{skel}}$.}
\label{tab:gran}
\end{table}

\subsection{Repeat-unit mass from SMILES}

All masses are computed from the SMILES with implicit hydrogens made explicit,
and with one correction that is easy to get wrong. Stripping the \texttt{[*]}
markers leaves a \emph{closed-shell} molecule: \texttt{[*]CC[*]} becomes
\texttt{CC}, which a cheminformatics toolkit dutifully caps with hydrogen,
giving ethane C$_2$H$_6$. But in a polymer those two carbons bond to
neighbouring repeat units, not to hydrogen. The repeat unit is C$_2$H$_4$.
PAAF therefore subtracts one hydrogen per attachment point,
\begin{equation}
  M_{\text{unit}} = \sum_{e} n_e\, m_e \;-\; n_{\ast}\, m_{\mathrm{H}} .
  \label{eq:mass}
\end{equation}
Without this correction every mass is high by $1.008$ per connection ---
$2.016$~g/mol for a linear unit, a $7\%$ error for polyethylene --- which
would propagate directly into cell density and into every weight fraction.

\subsection{Growth order}

Chains enter the box one at a time, so whichever species is grown \emph{last}
faces a box already at its final density. Taking species in the order the user
happens to list them therefore makes success depend on an arbitrary choice: a
$70/30$ PE/PS blend fails on the very first polystyrene chain, because 27
polyethylene chains have already filled the box and a polystyrene bead has
radius $2.74$~\AA{} against polyethylene's $1.77$.

PAAF therefore chooses the order by difficulty,
\begin{equation}
  \text{rank} \;=\; n_{\text{beads}} \times R_{\text{bead}}^{3},
  \qquad \text{descending},
  \label{eq:order}
\end{equation}
so the hardest species claims the empty box and the easiest fills the gaps.
The chosen order is reported in the result whenever it differs from the one
the caller supplied. Atoms are \emph{assembled} in the
caller's original order regardless, so that downstream consumers which slice
the atom list by species are unaffected.

\begin{caveat}
This reordering is an engineering addition by PAAF, not part of
Theodorou--Suter or Meirovitch, and is identified as such.
\end{caveat}

%======================================================================
\section{Composition: from formulation to cell}
\label{sec:composition}

A formulation is specified in weight percent; the grower needs integer chain
counts and a box edge. PAAF solves this explicitly rather than implicitly.

Given target weight fractions $w_i$ and chain masses
$m_i = M_{\text{unit},i}\,\mathrm{DP}_i$, the ideal fractional chain count is
\begin{equation}
  n_i = \frac{w_i\, M_{\text{total}}}{m_i} ,
  \label{eq:counts}
\end{equation}
and the integers are chosen by largest remainder: floor everything, then
allocate the leftover chains to whichever species is furthest below its
target. Every species requested with non-zero weight receives at least one
chain --- a ``5 wt\% additive'' that rounds to zero chains is not a 5 wt\%
cell, and silently dropping it would be the worst available outcome.

Dividing by $m_i$ in Eq.~\eqref{eq:counts} is essential and is the step most
often got wrong. A $\mathrm{DP}=50$ polystyrene chain is $5\,208$~g/mol
against $1\,403$ for polyethylene. Allocating chains in proportion to weight
percent directly would give $7:3$ for a $70/30$ blend and a realised
composition of $38.6$~wt\% PE. The correct solution is $18:2$, giving
$70.80$~wt\%.

The box edge follows from the target density $\rho$:
\begin{equation}
  a = \left(\frac{M_{\text{total}}}{N_A\,\rho}\right)^{1/3} \times 10^{8}\ \text{\AA} .
  \label{eq:box}
\end{equation}

Because chain counts are integers, the realised composition is generally not
the requested one. PAAF reports both, side by side, and states the error in
percentage points rather than rounding silently. The error falls with cell
size as Table~\ref{tab:comp} shows.

\begin{table}[t]
\centering
\small
\begin{tabular}{rrrr}
\toprule
Target beads & Chains & Realised wt\% PE & Error (pp) \\
\midrule
400    & 5   & 51.86 & 18.14 \\
2\,000  & 20  & 70.80 & 0.80 \\
10\,000 & 100 & 70.80 & 0.80 \\
40\,000 & 400 & 70.22 & 0.22 \\
\bottomrule
\end{tabular}
\caption{Composition error against cell size for a requested 70/30 PE/PS blend
at $\mathrm{DP}=50$. The error is a property of integers, not a defect;
PAAF states it rather than concealing it.}
\label{tab:comp}
\end{table}

%======================================================================
\section{From skeletal beads to atoms}
\label{sec:backmap}

Growth produces skeletal atoms. No force-field typer can work with those: DL\_FIELD
and Moltemplate need elements, hydrogens, side groups and bond orders. PAAF
therefore back-maps, and it can do so cleanly \emph{because} a bead is a
skeletal atom --- the backbone of the atomistic chain already exists, exactly,
at positions the growth accepted.

For each repeat unit PAAF embeds one 3D template: the repeat unit with its two
attachment points converted to hydrogens, minimised with MMFF (falling back
to UFF where MMFF typing fails). In that template
each backbone atom has a local frame defined by its two backbone neighbours
(a link hydrogen standing in at the unit ends). The same frame is built at the
corresponding bead in the grown chain, and the rigid transform between them,
\begin{equation}
  \mathbf{R} = \mathbf{F}_{\text{target}}\,\mathbf{F}_{\text{template}}^{\mathsf{T}},
  \qquad
  \mathbf{x}'_k = \mathbf{p}_{\text{bead}} + \mathbf{R}\,
  (\mathbf{x}_k - \mathbf{x}_{\text{backbone}}),
  \label{eq:backmap}
\end{equation}
is applied to everything attached to that backbone atom. Because
Eq.~\eqref{eq:backmap} is rigid, every substituent bond length and bond angle
survives exactly.

Verification: back-mapped formulae are exact --- a $\mathrm{DP}=n$
polyethylene chain gives C$_{2n}$H$_{4n+2}$ and a polystyrene chain
C$_{8n}$H$_{8n+2}$, confirmed for multi-chain cells --- and the backbone atoms
land on the grown beads to within $1.8\times10^{-15}$~\AA.

\paragraph{Push-off.} A hydrogen sits $1.09$~\AA{} from its carbon, so two
backbone atoms at an acceptable separation can still leave their hydrogens
nearly coincident. Handing a $0.3$~\AA{} contact to LAMMPS is how a
minimisation explodes. PAAF therefore applies a purely geometric relief step
in which \emph{only non-backbone atoms move}. It has two parts, and the order
matters.

First, each side group is \textbf{rotated rigidly} about the single bond
joining it to the backbone. That rotation is exact --- every bond length, bond
angle and internal torsion of the group is preserved, because the group moves
as a body --- and it is the \emph{only} move that can help a ring. It is
precisely the side-group torsion that growth does not sample, recovered here
at negligible cost. Each rotatable group is scanned over 36 uniform angles and
placed where a soft repulsive score against everything outside the group is
smallest.

Second, the remaining atoms are pushed apart, after which a SHAKE-style
constraint relaxation restores \textbf{every} bond, including ring closures,
to its original length.

Measured: polyethylene $0.34 \rightarrow 1.65$~\AA{} (no rotatable groups);
polystyrene $0.31 \rightarrow 1.00$~\AA{} with 20 phenyl groups rotated. Bond
lengths after both steps span $1.087$--$1.530$~\AA{} for polystyrene and
$1.094$--$1.530$~\AA{} for polyethylene, i.e.\ unchanged. Bond angles and
inter-group torsions do drift; that is the price, and it is the cheap part of
the geometry for a real minimiser to repair.

\begin{caveat}
\textbf{Bug history, recorded because the failure mode is instructive.}
An earlier implementation pinned each mobile atom to a \emph{single} parent
and re-projected onto that one bond. A spanning tree cannot constrain a cycle,
so the ring-closing bond of every phenyl group was left completely free and
was torn open to as much as $5.0$~\AA{} --- producing a structure that looked
plausible and was chemically destroyed. Nothing caught it, because the only
bond-length test in the suite used polyethylene, where every mobile atom is a
hydrogen bonded directly to a pinned backbone carbon and a tree is therefore
sufficient. \emph{Test the case that contains the cycle.}
\end{caveat}

\paragraph{Spearing, checked here.} Once the rings exist, the spearing test of
Eq.~\eqref{eq:spear} can finally be applied. Rings are located by a
depth-limited search for the shortest alternative path across each bond, then
approximated as discs by singular-value decomposition of their atom positions.
Every bond in the cell is tested against every ring under minimum image. A
polystyrene cell at $1.04$~g/cm$^3$ contained 20 phenyl rings and \textbf{3
speared bonds}; a polyethylene cell contains no rings and reports none. The
count is attached to the result and raised as a warning.

\paragraph{Tacticity.} Applying one template to every unit produces an
\emph{isotactic} chain. This is a real stereochemical choice that changes
properties, so PAAF exposes it rather than hiding it: \texttt{syndiotactic}
alternates and \texttt{atactic} mirrors each unit's substituents through the
local backbone plane at random, inverting the configuration at the
stereocentre.

%======================================================================
\section{The algorithm, in full}
\label{sec:algorithm}

\begin{algobox}
\algoreset
\small
\alg{\kw{input:} repeat-unit SMILES, DP and wt\% per species; target density $\rho$; temperature $T$}
\alg{compute $M_{\text{unit}}$ (Eq.~\ref{eq:mass}) and $n_{\text{skel}}$ (Eq.~\ref{eq:granularity}) for each species}
\alg{solve integer chain counts (Eq.~\ref{eq:counts}) and cubic box edge $a$ (Eq.~\ref{eq:box})}
\alg{order species by difficulty (Eq.~\ref{eq:order})}
\alg{\kw{for each} species $i$ in difficulty order:}
\alg[1]{\kw{for each} chain of species $i$:}
\alg[2]{\kw{for} attempt $=1 \dots N_{\text{restart}}$:}
\alg[3]{seed three atoms at a random position having local free volume}
\alg[3]{\kw{while} chain incomplete:}
\alg[4]{\kw{for each} rotational isomeric state $s \in \{t, g^{+}, g^{-}\}$:}
\alg[5]{place trial atom by NeRF (Eq.~\ref{eq:nerf})}
\alg[5]{\kw{reject} if overlap (Eq.~\ref{eq:overlap}) or ring spearing (Eq.~\ref{eq:spear})}
\alg[5]{evaluate $E^{\text{nb}}_{s}$ (Eq.~\ref{eq:softcore}); optionally look ahead (Eq.~\ref{eq:scan})}
\alg[5]{$w_s \leftarrow U[j,s]\,\exp(-E^{\text{nb}}_{s}/k_BT)$}
\alg[4]{\kw{if} all $w_s = 0$: \kw{restart chain} \quad\textit{(dead end --- attrition)}}
\alg[4]{\kw{else}: sample $s$ from $w_s / \sum w$; accept the atom}
\alg[2]{wrap the chain, insert into the neighbour cell list, record $R_g$, $R_{ee}$, $C_n$}
\alg{assemble atoms in the \emph{caller's} species order, not the growth order}
\alg{back-map to all-atom (Eq.~\ref{eq:backmap}); push off side-group overlaps}
\alg{assign atom types; export LAMMPS via DL\_FIELD or Moltemplate}
\end{algobox}
\vspace{4pt}
{\small\textbf{Algorithm 1:} PAAF amorphous cell construction.}

%======================================================================
\section{Validation}
\label{sec:validation}

Every physical claim PAAF makes is checked against an independent reference,
and the measured value is printed rather than merely asserted. The suite
comprises 442 tests, most of them a sweep of the whole polymer library; the GUI subset additionally requires PyQt5.

\subsection{Chain dimensions against literature}

Table~\ref{tab:cinf} is the load-bearing result. Deviations fall in both
directions, which is what one expects from a genuine model rather than a
fitted fudge factor.

\begin{table}[t]
\centering
\small
\begin{tabular}{lrrrrrrl}
\toprule
Polymer & $E_\sigma$ & $E_\omega$ & $\theta$ & $C_\infty$ meas. & Ref. & Dev. & Source of parameters \\
\midrule
PE   & 0.50 & 2.10 & $112^\circ$ & $6.26 \pm 0.25$ & 6.7  & $-6.6\%$  & Flory 1969~\cite{flory1969} \\
PP   & 0.20 & 1.65 & $114^\circ$ & $5.83 \pm 0.22$ & 5.5  & $+5.9\%$  & Suter \& Flory 1975~\cite{suter1975} \\
PS   & 0.85 & 2.40 & $114^\circ$ & $8.90 \pm 0.38$ & 10.0 & $-11.0\%$ & Yoon \emph{et al.} 1975~\cite{yoon1975} \\
PMMA & 0.70 & 2.20 & $116^\circ$ & $9.16 \pm 0.38$ & 8.2  & $+11.8\%$ & Sundararajan \& Flory 1974~\cite{sundararajan1974} \\
\bottomrule
\end{tabular}
\caption{Characteristic ratio against tabulated values. Energies in kcal/mol,
$T = 413$~K, 400 chains of 500 bonds each. The PP, PS and PMMA entries average
over tacticity rather than modelling the meso/racemo distinction explicitly,
which is the principal source of their deviation.}
\label{tab:cinf}
\end{table}

\subsection{Convergence and internal consistency}

$C_n$ must increase with chain length and level off, not sit flat.
Measured for PE at 413~K: $5.01$ ($n=20$), $6.33$ ($50$), $6.30$ ($100$),
$7.04$ ($200$), $7.21$ ($400$). Cooling must favour trans, and does:
the trans fraction rises $0.573 \rightarrow 0.602 \rightarrow 0.643$ as
$T$ falls $500 \rightarrow 400 \rightarrow 300$~K.

\subsection{Coverage across the polymer library}

Validating two polymers demonstrates that two polymers work. Every entry in
PAAF's 124-polymer library is therefore driven through chemistry, growth and
back-mapping, and each must either yield sound bond lengths or be refused with
a specific reason. The sweep found real defects that PE and PS could not have
exposed:

\begin{itemize}
\item \textbf{Poly(methyl methacrylate) could not be parsed at all.} Its
      repeat unit \texttt{[*]CC([*])(C)C(=O)OC} has a quaternary backbone
      carbon, so the tail after the second marker \emph{already} begins with a
      branch. Wrapping it in another pair of parentheses produced
      \texttt{CC((C)C(=O)OC)} --- a doubled bracket that is not valid SMILES.
      One of the four polymers whose RIS parameters this report tabulates
      could not have its mass computed. Fixed by peeling the leading balanced
      branches before parenthesising the remainder; PMMA now gives
      $100.12$~g/mol (C$_5$H$_8$O$_2$), exact.
\item \textbf{Poly(dimethylsiloxane) crashed.} Both attachment points sit on
      the same silicon, and RDKit raises an invariant violation when asked for
      the shortest path from an atom to itself. The single-skeletal-atom case
      is now handled directly; PDMS gives Si$_{16}$C$_{32}$H$_{100}$ for
      2 chains of $\mathrm{DP}=8$, exact.
\item \textbf{24 polymers have a ring in the backbone} and were silently
      producing torn structures with bonds up to $11.9$~\AA. They are now
      refused with a specific explanation (limitation L9).
\item \textbf{One library entry, MAH, carries a non-polymerisable SMILES} in
      its polymerisation field. This is a data defect rather than a code
      defect and is asserted against so it cannot spread unnoticed.
\end{itemize}

\subsection{Cell-level checks}

\begin{itemize}
\item \textbf{Density.} Target densities are reached exactly: $0.8500$,
      $0.8800$, $0.9000$, $0.9500$, $1.0400$~g/cm$^3$ all realised to four
      decimal places, because the box is sized analytically from
      Eq.~\eqref{eq:box}. Of these, $0.85$ and $1.04$ are reached with
      unbiased sampling; $0.88$, $0.90$ and $0.95$ come from tests that
      enable $k=1$.
\item \textbf{Geometry preserved under growth.} Bond length
      $1.530000 \pm 1.4\times10^{-15}$~\AA, angle
      $112.0000 \pm 5.1\times10^{-14}\,^\circ$ in a packed cell.
\item \textbf{No overlaps.} Minimum non-bonded distance under minimum image
      measured at $2.002$~\AA{} for a requested tolerance of $2.0$~\AA.
\item \textbf{Spearing predicate.} Unit-tested against four geometric cases:
      straight through, passing beside, parallel to the ring plane, and
      stopping short.
\item \textbf{Reproducibility.} Identical seed gives identical coordinates to
      $0.00\times10^{0}$~\AA{} maximum difference.
\item \textbf{Melt-level $C_n$.} Polystyrene in a cell at $0.95$~g/cm$^3$
      gives $11.13$ against a reference $C_\infty$ of $10.0$.
\end{itemize}

%======================================================================
\section{Limitations}
\label{sec:limitations}

This section exists because a reader will ask these questions, and answering
them first is stronger than being caught by them.

\begin{enumerate}[label=\textbf{L\arabic*.}]

\item \textbf{Construction only --- there is no energy minimisation.}
BIOVIA describes Amorphous Cell as two steps: an initial guess structure,
\emph{followed by relaxation to a state of minimum potential
energy}~\cite{biovia}. PAAF performs only the first. The geometric push-off of
Section~\ref{sec:backmap} makes the structure safe to hand to a minimiser; it
is not itself a minimisation and involves no force field. \emph{Users must
minimise and equilibrate in LAMMPS before quoting any property from a PAAF
cell.} This is the single largest gap between PAAF and a production tool.

\item \textbf{Growth is at united-skeletal-atom resolution.}
Theodorou and Suter place every pendant atom and sample side-group torsions
for vinyl polymers. PAAF grows the backbone and reconstructs side groups
afterwards from a single template rotamer, then rotates each group about its
own bond to relieve overlap. That rotation recovers one degree of freedom
exactly, but it is a greedy clash-relief step, \emph{not} Boltzmann sampling
of the side-group torsion: the resulting rotamer distribution has no claim to
be the equilibrium one. Properties that depend on side-group packing should
not be taken from an unequilibrated PAAF cell.

\item \textbf{Spearing is detected but cannot be prevented.}
Because rings do not exist until back-mapping, a speared ring can only be
reported after the fact, not rejected during growth. This is not hypothetical:
a polystyrene cell at $1.04$~g/cm$^3$ contained 3 speared bonds among 20
rings. PAAF raises a warning; the user must rebuild with a different seed or
at lower density. A full-atomistic grower would reject these as they arise.

\item \textbf{The non-bonded term is generic, not the production force field.}
Equation~\eqref{eq:softcore} is a soft repulsion, not the Lennard-Jones and
electrostatic terms of the force field that will later be applied. Materials
Studio uses the real force field during growth.

\item \textbf{RIS parameters exist for four polymers.}
PE, PP, PS and PMMA carry literature parameterisations. Every other polymer
falls back to a generic three-state model, which PAAF labels
\texttt{is\_generic} and flags in the result as ``chain dimensions are
qualitative, not a quantitative prediction''. It is honest, but it means the
RIS advantage does not extend to arbitrary chemistry. This is the most
consequential limitation for a user working outside the four.

\item \textbf{Tacticity is averaged in the RIS parameters.}
PP, PS and PMMA use a single averaged state set in place of the meso/racemo
distinction of the published treatments, which is the main reason PMMA reads
$+11.8\%$.

\item \textbf{Scanning is biased and therefore off by default.}
See Section~\ref{sec:growth}. The Rosenbluth weight is not carried, so
$k>0$ cells must not be used to quote chain dimensions.

\item \textbf{Sequential construction over-extends chains by roughly 20\%.}
Chains grow into a box that starts empty, so later chains encounter free
volume more open and more connected than a real melt's, and
Eq.~\eqref{eq:growth} then favours the extended conformations that fit it.
Measured for PE at $0.85$~g/cm$^3$, $\mathrm{DP}=60$: $C_n = 8.79$ in the cell
against $7.26$ for the same chain grown in isolation. Real melts screen
excluded volume and sit near unperturbed dimensions, so a constructed cell's
$C_n$ should be treated as an upper bound and re-measured after equilibration.

\item \textbf{Linear chains only.} No branching, no polydispersity, no
dissolved small molecules, no layered or oriented phases --- all of which
Materials Studio supports~\cite{biovia,accelrys}.

\item \textbf{Backbone rings are refused, not supported.} Back-mapping assumes
the groups hanging off the backbone are \emph{trees}. A ring \emph{in} the
backbone --- poly(ethylene terephthalate), polycarbonate, PBAT --- is a cycle
spanning several skeletal atoms whose internal geometry is fixed by the ring,
not by the RIS bond length and angle the growth used. The two are
incompatible. \textbf{24 of the 123 usable library polymers} are of this kind.
PAAF now detects them and raises a specific error; growth and coarse-grained
export still work, since a bead is a skeletal atom either way, but no all-atom
cell and therefore no force field is possible. Pendant rings (polystyrene) are
unaffected.

\item \textbf{The cell must be larger than twice the longest repeat unit.}
Minimum image requires it. A repeat unit with a long side chain --- an
octyloxy group is $11.1$~\AA{} across --- violates it in a small cell, and
bonds measured across the boundary are then wrapped to the wrong periodic
image. Measured for \texttt{[*]C(OCCCCCCCC)C[*]}: bonds up to $12.29$~\AA{} at
2 chains of $\mathrm{DP}=6$ (box $15.4$~\AA), and $1.093$--$1.532$~\AA{} for
the same polymer at 4 chains of $\mathrm{DP}=20$ (box $29.0$~\AA). The
structure was never wrong; the cell was too small to describe it. PAAF now
compares the template extent against half the box and warns.

\item \textbf{Some library SMILES layouts are not expanded reliably.}
The $n$-mer expansion recognises the two common layouts, \texttt{[*]A[*]} and
\texttt{[*]A([*])B}. Entries that put a marker inside a branch, use a bare
\texttt{*}, or do not begin with a marker fall back to stripping the markers
and repeating, which is not always correct, and log a warning when they do.
Repeat-unit \emph{masses} are unaffected because they need only one unit.
\end{enumerate}

%======================================================================
\section{Relation to existing tools, and an honest novelty assessment}
\label{sec:novelty}

\subsection{Comparison}

\begin{table}[h]
\centering
\small
\begin{tabular}{p{4.3cm}p{4.6cm}p{5.2cm}}
\toprule
 & \textbf{Materials Studio} & \textbf{PAAF} \\
\midrule
Growth method    & Theodorou--Suter + Meirovitch & same \\
Chain statistics & RIS & RIS, validated against literature $C_\infty$ \\
Resolution       & full atomistic & united skeletal atoms, then back-mapped \\
Non-bonded during growth & production force field & generic soft repulsion \\
After construction & relaxes to potential-energy minimum & \emph{nothing} \\
System types     & polydisperse, branched, gases, layers, nematics & linear homopolymers and blends \\
Licence          & proprietary & open, scriptable, seed-reproducible \\
\bottomrule
\end{tabular}
\caption{PAAF against BIOVIA Materials Studio Amorphous Cell~\cite{biovia,accelrys}.}
\end{table}

Among open-source tools, RadonPy~\cite{radonpy} automates SMILES $\rightarrow$
amorphous cell $\rightarrow$ MD $\rightarrow$ properties, pysimm~\cite{pysimm}
provides polymer construction into LAMMPS, and SPACIER~\cite{spacier} builds
on RadonPy for machine-learning-driven design. These construct cells
predominantly by self-avoiding random walk with packing and compression,
relying on extended equilibration to relax chain dimensions, rather than by
RIS-weighted growth.

\subsection{Novelty}

Stated plainly: \textbf{the method implemented here is not novel.} It is a
faithful reimplementation of Theodorou--Suter (1985) and Meirovitch (1983) on
Flory's RIS foundation (1969). Nothing in the growth algorithm would survive a
novelty assessment. The composition solver, the difficulty ordering and the
back-mapping are competent engineering, not new science.

What is arguably underserved is the combination: \emph{RIS-correct
construction in an open-source tool}. Materials Studio has it and is closed;
the widely used open tools largely do not. That suggests one testable claim,
which would need to be demonstrated rather than asserted:

\begin{quote}
A cell built with RIS statistics begins nearer the correct chain dimensions
and therefore reaches equilibrium in less MD time than a random-walk cell of
identical composition.
\end{quote}

Testing it requires building matched systems by both routes, running NPT
equilibration, and tracking $C_\infty$ and $R_g$ against simulation time. If
the crossover is substantial, that is a modest and publishable software
result. If it is not, the tool's value is convenience and transparency rather
than physics --- which is worth establishing before making any claim.

Two developments would constitute genuine contributions. The first, and
stronger: \textbf{deriving RIS parameters automatically from SMILES} by
torsion scanning a trimer fragment at MMFF, GFN2-xTB or DFT level and fitting
$E_\sigma$, $E_\omega$ and the state angles. This would extend RIS-correct
construction from four polymers to arbitrary chemistry and would remove
limitation~L4, which is the one that most restricts the method's usefulness.
The second: carrying the Rosenbluth weight through scanning so that dense
cells and unbiased chain dimensions can be obtained simultaneously, removing
limitation~L6.

%======================================================================
\begin{thebibliography}{99}
\small

\bibitem{ts1985}
D.~N. Theodorou and U.~W. Suter,
\emph{Detailed molecular structure of a vinyl polymer glass},
Macromolecules \textbf{18} (1985) 1467--1478.

\bibitem{meirovitch1983}
H.~Meirovitch,
\emph{Computer simulation of self-avoiding walks: Testing the scanning method},
J. Chem. Phys. \textbf{79} (1983) 502--508.

\bibitem{flory1969}
P.~J. Flory,
\emph{Statistical Mechanics of Chain Molecules},
Interscience, New York, 1969.

\bibitem{mattice1994}
W.~L. Mattice and U.~W. Suter,
\emph{Conformational Theory of Large Molecules: The Rotational Isomeric State
Model in Macromolecular Systems},
Wiley, New York, 1994.

\bibitem{suter1975}
U.~W. Suter and P.~J. Flory,
\emph{Conformational energy and configurational statistics of polypropylene},
Macromolecules \textbf{8} (1975) 765--776.

\bibitem{yoon1975}
D.~Y. Yoon, P.~R. Sundararajan and P.~J. Flory,
\emph{Conformational characteristics of polystyrene},
Macromolecules \textbf{8} (1975) 776--783.

\bibitem{sundararajan1974}
P.~R. Sundararajan and P.~J. Flory,
\emph{Configurational characteristics of poly(methyl methacrylate)},
J. Am. Chem. Soc. \textbf{96} (1974) 5025--5031.

\bibitem{rosenbluth1955}
M.~N. Rosenbluth and A.~W. Rosenbluth,
\emph{Monte Carlo calculation of the average extension of molecular chains},
J. Chem. Phys. \textbf{23} (1955) 356--359.

\bibitem{grassberger1997}
P.~Grassberger,
\emph{Pruned-enriched Rosenbluth method: Simulations of $\theta$ polymers of
chain length up to 1\,000\,000},
Phys. Rev. E \textbf{56} (1997) 3682--3693.

\bibitem{biovia}
BIOVIA,
\emph{What Does BIOVIA Materials Studio Amorphous Cell Do?},
Dassault Syst\`emes.
\url{https://www.3ds.com/assets/invest/2023-10/biovia-material-studio-amorphous-cell.pdf}

\bibitem{accelrys}
Accelrys,
\emph{Amorphous\_Cell Module}, Polymer User Guide.
\url{http://www.esi.umontreal.ca/accelrys/materials/insight400P/polymer/03-AC.html}

\bibitem{radonpy}
Y.~Hayashi \emph{et al.},
\emph{RadonPy: automated physical property calculation using all-atom
classical molecular dynamics simulations for polymer informatics},
npj Comput. Mater. \textbf{8} (2022) 222.

\bibitem{pysimm}
M.~E. Fortunato and C.~M. Colina,
\emph{pysimm: A python package for simulation of molecular systems},
SoftwareX \textbf{6} (2017) 7--12.

\bibitem{spacier}
H.~Arai \emph{et al.},
\emph{SPACIER: on-demand polymer design with fully automated all-atom
classical molecular dynamics integrated into machine learning pipelines},
npj Comput. Mater. \textbf{10} (2024).

\bibitem{parsons1994}
D.~F. Parsons, D.~R.~M. Williams \emph{et al.};
see also J.~Gasteiger and T.~Engel (eds.),
\emph{Chemoinformatics}, Wiley-VCH, 2003, for the NeRF internal-coordinate
construction used in Eq.~\eqref{eq:nerf}.

\end{thebibliography}

\end{document}
